\subsection{Vergleich der Messwerte}
\begin{table}[htb]
  \centering
  \caption{Vergleich der Messergebnisse für \(\kappa\) mit den Literaturwerten von \(\kappa_\text{L,lit} = 1.401\) und \(\kappa_\text{A,lit} = 1.648\)}
  \begin{tabularx}{0.8\textwidth}{TTZZ}
    \toprule
         Medium  & Verfahren             & \kappa            & S\,[\sigma]\\
        \midrule
        Luft    & Clément-Desormes     & \num{1.377(14)}   & \num{1.8}\\
                & Rüchardt              & \num{1.357(18)}   &\num{2.5}\\
         \midrule
        Argon   & Rüchardt              & \num{1.603(23)}& \num{2.0}\\
    \bottomrule
  \end{tabularx}
\label{table:kappa_vergleich}  
\end{table}
Angenehme Signifikanzen. 
\begin{table}[htb]
  \centering
  \caption{Vergleich des \(\kappa_\text{Luft}\) Wertes nach Clément-Desormes mit Rüchardt}
  \begin{tabularx}{0.8\textwidth}{TZZZ}
    \toprule
        Medium  & \kappa^\text{CD}   & \kappa^\text{R}  & S\,[\sigma]\\
        \midrule
        Luft    &  \num{1.377(14)}    & \num{1.357(18)}  &\num{0.15}\\
    \bottomrule
  \end{tabularx}
\label{table:kappa_Luft}  
\end{table}

Die erste Beobachtung stellt fest, dass die Signifikanzen im grünen Bereich liegen. 

Zweitens lässt sich beobachten, dass die Ergebnisse für alle Adiabatenkoeffizient systematisch kleiner sind als der Literaturwert. Die könnte an der Luftzusammensetzung liegen. So könnte der Anteil an Gasen mit kleinerem Adiabatenkoeffizienten höher sein, sodass die Näherung \(\kappa_\text{Luft}\approx\kappa_\mathrm{N_2}\) nicht gerechtfertigt ist.  

Eine positive Beobachtung ist, dass die Messergebnisse der zwei Labormethoden zueinander konsistent sind. Dies deutet darauf hin, dass ein systematischer Fehler beide Versuche beeinflusst hat, wie etwa der Luftdruck oder die Raumtemperatur.\\
Evlt. können Abweichungen zu Literaturwerten auch den vorgenommenen Näherungen in der Formelherleitung für Clément-Desormes geschuldet sein, oder ungenauem Ablesen der Höhenpegelstände. Zudem könnte die Genauigkeit der Rüchardt Methode durch Aufnehmen einer ganzen Messreihe an Periodendauern noch verbessert werden. 


\subsection{Fazit}
Die Adiabatenkoeffizienten konnten zufriedenstellen genähert werden, bis auf einen wahrscheinlich systematischen Fehler bei Druck, Luftfeuchtigkeit oder Temperatur im Raum, welche sich nach Aufschreiben der einmaligen Messwerte evtl. verändert haben.

Schöner schneller Versuch. Nur ist mir bisher kein Meme eingefallen.
